% ============================================================
%  Come Consciente — Guía práctica sobre aditivos alimentarios
%  faustobe.it | 2026 | CC BY-NC 4.0
% ============================================================
\documentclass[a5paper,12pt,oneside,openany]{book}

\usepackage[T1]{fontenc}
\usepackage[utf8]{inputenc}
\usepackage{ebgaramond}
\usepackage[a5paper, top=2cm, bottom=3cm, inner=2.2cm, outer=3cm]{geometry}
\usepackage[table]{xcolor}
\usepackage[spanish]{babel}

% Colors
\definecolor{verdescuro}{HTML}{2d6a4f}
\definecolor{verdecharo}{HTML}{52b788}
\definecolor{grigioscuro}{HTML}{2b2d42}
\definecolor{amarillo}{HTML}{d4a017}
\definecolor{rojo}{HTML}{c1121f}
\definecolor{azulinfo}{HTML}{3a86ff}

\usepackage{tcolorbox}
\tcbuselibrary{skins,breakable}

\usepackage{booktabs}
\usepackage{tabularx}
\newcolumntype{L}{>{\raggedright\arraybackslash}X}

\usepackage{hyperref}
\hypersetup{
  colorlinks=true,
  linkcolor=verdescuro,
  urlcolor=verdescuro,
  pdftitle={Come Consciente},
  pdfauthor={faustobe},
}

\usepackage{qrcode}
\usepackage{fancyhdr}
\usepackage{titlesec}
\usepackage{microtype}

% Colored boxes
\tcbset{
  cajabase/.style={
    breakable, enhanced, arc=4pt, boxrule=0.8pt,
    fonttitle=\bfseries\small,
    before upper={\parskip=4pt},
  }
}

\newcommand{\boxverde}[2]{%
  \begin{tcolorbox}[cajabase,
    colback=verdescuro!8!white, colframe=verdescuro, title={#1}]
  #2\end{tcolorbox}}

\newcommand{\boxamarillo}[2]{%
  \begin{tcolorbox}[cajabase,
    colback=amarillo!12!white, colframe=amarillo!80!black, title={#1}]
  #2\end{tcolorbox}}

\newcommand{\boxrojo}[2]{%
  \begin{tcolorbox}[cajabase,
    colback=rojo!8!white, colframe=rojo, title={#1}]
  #2\end{tcolorbox}}

\newcommand{\boxinfo}[2]{%
  \begin{tcolorbox}[cajabase,
    colback=azulinfo!8!white, colframe=azulinfo, title={#1}]
  #2\end{tcolorbox}}

% Page style
\pagestyle{fancy}
\fancyhf{}
\fancyfoot[C]{\footnotesize\textcolor{grigioscuro}{\textit{Come Consciente\ \textbar\ faustobe.it}}}
\fancyfoot[R]{\footnotesize\thepage}
\renewcommand{\headrulewidth}{0pt}
\renewcommand{\footrulewidth}{0.4pt}
\fancypagestyle{plain}{%
  \fancyhf{}%
  \fancyfoot[C]{\footnotesize\textcolor{grigioscuro}{\textit{Come Consciente\ \textbar\ faustobe.it}}}%
  \fancyfoot[R]{\footnotesize\thepage}%
  \renewcommand{\headrulewidth}{0pt}%
  \renewcommand{\footrulewidth}{0.4pt}%
}

% Chapter title style
\titleformat{\chapter}[display]
  {\normalfont\LARGE\bfseries\color{verdescuro}}
  {\chaptertitlename\ \thechapter}{10pt}{\Huge}
\titlespacing*{\chapter}{0pt}{-20pt}{30pt}

% ============================================================
\begin{document}

% -------- PORTADA --------
\begin{titlepage}
\pagecolor{verdescuro}
\color{white}
\vspace*{2.5cm}
\begin{center}
  {\fontsize{38}{44}\selectfont\bfseries Come Consciente}\\[1.2cm]
  \textcolor{verdecharo}{\rule{6cm}{0.6pt}}\\[1.2cm]
  {\large\itshape Guía práctica sobre aditivos alimentarios}\\[0.6cm]
  {\normalsize Con E-Codes Reader para Android}\\[3cm]
  {\normalsize\textcolor{verdecharo}{faustobe.it}}
\end{center}
\vfill
\begin{center}
  {\small Versión 1.0\ \textbullet\ 2026\ \textbullet\ CC BY-NC 4.0}
\end{center}
\end{titlepage}
\nopagecolor

\frontmatter
\tableofcontents
\mainmatter

% ============================================================
\chapter{Cómo usar la app para hacer la compra de forma consciente}

Aprende a reconocer los alimentos más procesados y elegir los verdaderamente simples
y naturales en pocos minutos, sin convertirte en un experto en aditivos.

\section{Antes de escanear, haz tres comprobaciones rápidas}

Antes de abrir la app, observa el envase. Las etiquetas frontales están a menudo
diseñadas para que los productos parezcan más saludables de lo que son.

\begin{itemize}
  \item \textbf{Mira el frente del envase.} Presta atención a reclamos como
    «sin azúcar pero con edulcorantes», «sin grasas pero con emulsionantes»,
    «sabor extra» o «aromatizado». Suelen indicar que el producto ha sido
    muy modificado.
  \item \textbf{Abre la app y escanea el producto.} Usa la cámara para escanear
    el código de barras o sube una foto de la etiqueta. La app te mostrará de
    inmediato la lista de ingredientes y los aditivos presentes.
  \item \textbf{Comprueba tres puntos clave.} ¿Cuántos aditivos aparecen
    (sobre todo los códigos~E)? ¿Qué tan larga y compleja es la lista de
    ingredientes? ¿Cuántos azúcares, sal y grasas se declaran?
\end{itemize}

Si respondes «muchos» o «larga/complicada» a más de una pregunta, probablemente
estás ante un alimento muy procesado.

\section{Tres reglas para detectar un alimento «demasiado procesado»}

\textbf{Pocos ingredientes = mejor.} Si la lista es corta y contiene nombres
conocidos (por ejemplo «harina, agua, sal, aceite»), el producto está poco
procesado y suele ser más fiable.

\textbf{Muchos ingredientes «impronunciables».} Términos como «almidones
modificados», «proteínas aisladas», «mono y diglicéridos de ácidos grasos»,
«potenciadores del sabor» o «aromas complejos» son una señal clara de un alimento
muy procesado.

\textbf{Muchos códigos~E y aditivos.} No todos los~E son peligrosos, pero su
presencia masiva ---colorantes, conservantes, estabilizantes, emulsionantes---
indica un alto grado de transformación industrial.

\section{Cómo interpretar el resultado de la app}

Cada producto recibe una valoración rápida: piensa en estos colores como un
semáforo alimentario.

\boxverde{Poco procesado --- luz verde}{Pocos ingredientes, pocos o ningún
aditivo~E. Azúcares y sal moderados. Excelente opción para la compra diaria.}

\boxamarillo{Moderadamente procesado --- con moderación}{Algunos~E, aromas
naturales o azúcares añadidos. Está bien de vez en cuando, no como base de
cada comida.}

\boxrojo{Ultraprocesado --- limita o evita}{Muchos aditivos, azúcares añadidos,
aromas complejos y grasas de origen industrial. Mejor reservarlo como excepción,
especialmente para los niños.}

\section{Cómo usar la app de forma práctica en el supermercado}

\textbf{1.~Empieza por la sección de refrigerados y desayuno.} Escanea yogures,
pan, leche, bebidas. Busca productos con pocos ingredientes, pocos~E, y azúcares
simples (fruta, leche, azúcar de caña) en lugar de jarabes y edulcorantes.

\textbf{2.~Luego pasa a los dulces y snacks.} Si la lista es larguísima y está
llena de~E, la app marcará en rojo. Elige siempre la alternativa más sencilla:
fruta, pan integral, frutos secos.

\textbf{3.~Usa la app para comparar productos similares.} Compara dos productos
del mismo tipo en la app. A menudo el que tiene menos ingredientes es solo un
poco menos llamativo en el estante, pero mucho más natural.

\textbf{4.~Deja los «rojos» como excepción.} No hace falta eliminar todo lo
ultraprocesado: basta con reducirlo. Fíjate un límite ---máximo 1--2 productos
«rojos» a la semana--- mientras el resto de la compra se base en alimentos sencillos.

\section{Tabla resumen}

{\small
\begin{tabularx}{\textwidth}{|L|L|}
\hline
\rowcolor{verdescuro}\textcolor{white}{\textbf{Resultado en la app}}
  & \textcolor{white}{\textbf{Qué hacer en la práctica}} \\
\hline
Pocos ingredientes, pocos/ningún~E
  & Base diaria: yogur natural, pan integral, tomate triturado \\
\hline
Algunos~E, azúcares moderados
  & De vez en cuando, no como alimento principal \\
\hline
Muchos~E, azúcares añadidos, aromas complejos
  & Limita o evita; úsalo solo como excepción \\
\hline
\end{tabularx}
}

\section{Por qué aprender a leer los aditivos marca la diferencia}

Reconocer los alimentos más procesados no es «fobia a los aditivos»: es atención
a la calidad de lo que comemos. Numerosos estudios sugieren que reducir los
ultraprocesados ayuda a largo plazo en el control del peso, la salud cardiovascular
y la estabilidad metabólica.

La app te ofrece una herramienta concreta: no hace falta memorizar cada código~E,
solo aprender a reconocer las señales de exceso de procesamiento.

% ============================================================
\chapter{Cómo leer la etiqueta nutricional}

Ingredientes, calorías, sal oculta: aprende a descifrar una etiqueta en menos de
un minuto y deja de comprar productos que parecen sanos pero no lo son.

\section{La lista de ingredientes: por dónde empezar}

Los ingredientes aparecen en \textbf{orden decreciente de peso}: el primero es el
más abundante en el producto.

\begin{itemize}
  \item Si azúcar o aceite de palma encabezan la lista, el producto se construye
    alrededor de ese componente.
  \item La longitud de la lista indica el grado de procesamiento: el pan casero
    tiene 4 ingredientes, el pan industrial puede tener 15 o más.
  \item Términos como «almidón modificado» o «carboximetilcelulosa» indican
    procesamiento industrial.
\end{itemize}

\boxinfo{Regla práctica}{Si no encontrarías todos los ingredientes en un
supermercado normal, el producto es probablemente muy procesado.}

\section{La tabla nutricional: qué mirar (y qué ignorar)}

\begin{description}
  \item[Azúcares] Menos de 5\,g por 100\,g es bajo; entre 5 y 22,5\,g es medio;
    por encima de 22,5\,g es alto. Para bebidas, el límite bajo es 2,5\,g/100\,ml.
  \item[Sal] Menos de 0,3\,g/100\,g es bajo; más de 1,5\,g/100\,g es alto.
    El sodio debe multiplicarse por 2,5 para obtener la sal.
  \item[Grasas saturadas] Menos de 1,5\,g/100\,g es aceptable; más de 5\,g/100\,g
    se considera alto.
\end{description}

\section{Raciones vs 100\,g: el truco que confunde}

Las empresas destacan valores por ración que parecen menores, pero suelen estar
subestimados. Para comparar productos, usa siempre los valores por 100\,g. Cuando
una bolsa contiene múltiples raciones, multiplica los valores por el número total.

\boxamarillo{Ejemplo}{Si el envase indica 10\,g de azúcares por porción de 30\,g
y tú comes 60\,g, consumes 20\,g de azúcares ---casi la mitad de la
recomendación diaria.}

\section{La sal oculta: dónde se esconde de verdad}

El sal aparece también en productos aparentemente dulces.

\begin{itemize}
  \item \textbf{Pan y productos de panadería:} dos rebanadas pueden contener
    0,5--0,8\,g de sal.
  \item \textbf{Cereales de desayuno:} incluso versiones «integrales» contienen
    0,8--1,2\,g de sal por 100\,g.
  \item \textbf{Salsas y condimentos:} pueden alcanzar 1--3\,g de sal por 100\,g.
\end{itemize}

\section{Las declaraciones nutricionales: qué significan de verdad}

\textbf{«Sin azúcares añadidos»:} no significa ausencia total de azúcares;
puede contener fructosa o lactosa naturales y edulcorantes artificiales.

\textbf{«Light» o «reducido»:} indica una reducción del 30\,\% respecto al
producto de referencia, pero la grasa eliminada frecuentemente se reemplaza
con azúcares o espesantes.

\textbf{«Rico en fibra» o «Fuente de proteínas»:} cumplen umbrales legales
mínimos pero no aseguran ausencia de otros ingredientes problemáticos.

\section{Lista de verificación rápida para el supermercado}

\begin{enumerate}
  \item Examina el primer ingrediente: ¿es un alimento real o un azúcar/grasa
    procesados?
  \item Cuenta ingredientes: menos de 5 es excelente; 5--10 es aceptable; más
    de 10--12 requiere lectura atenta.
  \item Verifica los códigos~E: pocos o ninguno es positivo.
  \item Revisa sal y azúcares por 100\,g usando los umbrales indicados.
  \item Ignora las declaraciones en la cara frontal; confía en la lista de
    ingredientes.
\end{enumerate}

\boxverde{El truco más sencillo}{Si reconoces todos los ingredientes y podrías
cocinarlos en casa, vas bien encaminado.}

% ============================================================
\chapter{Aditivos frente a ultraprocesados: ¿cuál es la diferencia?}

¿Un producto con muchos aditivos es necesariamente malo? ¿Y un ultraprocesado
siempre está lleno de aditivos? La respuesta te sorprenderá.

\section{Qué son los aditivos alimentarios}

Los aditivos son sustancias añadidas intencionalmente a los alimentos para cumplir
funciones tecnológicas: conservar, colorear, espesar, emulsionar, estabilizar o
mejorar el sabor. En Europa, cada aditivo aprobado recibe un código~E y debe pasar
una evaluación de seguridad de la EFSA.

\begin{itemize}
  \item No todos los aditivos son sintéticos: el ácido cítrico, la curcumina
    y la lecitina de soja son de origen natural.
  \item La mayoría de los~E aprobados se considera segura en las dosis permitidas.
  \item El riesgo está en la \emph{cantidad} y la \emph{combinación} de múltiples
    aditivos.
  \item Algunos merecen especial atención: colorantes azoicos, conservantes de
    nitrato y edulcorantes intensivos.
\end{itemize}

\section{Qué son los alimentos ultraprocesados}

El término «ultraprocesado» proviene de la clasificación NOVA de la Universidad
de São Paulo, que se enfoca en los procesos industriales, no en los nutrientes.

Un producto NOVA~4 contiene ingredientes ausentes en cocinas domésticas: proteínas
aisladas, almidones modificados, hidrolizados proteicos, jarabes de
glucosa-fructosa, aromas «naturidénticos», emulsionantes, colorantes y
estabilizantes.

\textbf{Ultraprocesados:} bollería, snacks envasados, cereales azucarados,
bebidas gaseosas, embutidos reconstituidos, platos congelados con salsas,
yogures aromatizados, pan de molde industrial.

\textbf{No ultraprocesados:} quesos curados, charcutería de calidad, pan artesanal,
pasta seca tradicional, conservas de tomate sin aditivos, vinagre, aceite
de oliva virgen extra.

\section{Cuándo coinciden --- y cuándo no}

\begin{description}
  \item[Con aditivos pero NO ultraprocesado:] queso con nitrato potásico E252,
    vino con sulfitos E220, aceite con vitamina~E E306 como antioxidante.
  \item[Ultraprocesado con POCOS aditivos~E:] pan de molde con almidón
    modificado y aromas, barritas proteicas con proteínas aisladas, surimi.
  \item[Ultraprocesado Y con muchos aditivos:] bollería industrial, salsas
    industriales, bebidas energéticas.
\end{description}

\boxinfo{Conclusión}{Un producto con 15 ingredientes industriales sin ningún~E
sigue siendo un ultraprocesado. Fíjate en las dos cosas.}

\section{El doble problema: por qué ambos importan}

Los riesgos para la salud provienen de dos direcciones:

\textbf{Riesgos por aditivos específicos:} los colorantes azoicos se vinculan
con hiperactividad infantil; los nitratos en carnes procesadas forman
nitrosaminas potencialmente cancerígenas; los edulcorantes intensivos pueden
alterar el microbioma.

\textbf{Riesgos por el ultraprocesamiento en sí:} estudios a gran escala
(NutriNet-Santé, UK Biobank) asocian un consumo elevado de NOVA~4 con mayor
riesgo de obesidad, diabetes tipo~2, enfermedades cardiovasculares y ciertos
cánceres.

\textbf{El efecto combinado:} en la vida real consumimos combinaciones de
10--20 aditivos distintos durante la misma comida. Los efectos sinérgicos
aún están poco investigados.

\section{Cómo reconocerlos en la práctica}

\begin{enumerate}
  \item Identifica ingredientes «de laboratorio»: almidones modificados,
    proteínas aisladas, jarabes de glucosa-fructosa.
  \item Cuenta los códigos~E y clasifícalos en las categorías más críticas.
  \item Ten en cuenta el contexto dietético completo.
  \item Usa la app como ayuda, no como árbitro absoluto.
\end{enumerate}

% ============================================================
\chapter{Clasificación NOVA: qué es y cómo usarla}

NOVA no mide los nutrientes sino el \emph{grado de transformación industrial}
del alimento. Entender los 4 grupos cambia la forma en que haces la compra.

\section{Qué es la clasificación NOVA}

NOVA fue creado por el profesor Carlos Monteiro de la Universidad de São Paulo
y es adoptado por la OMS y la FAO. A diferencia del Nutri-Score, que evalúa la
composición nutricional, NOVA clasifica según el \emph{grado de transformación
industrial}: no el número de calorías, sino cuánta intervención industrial ha
sufrido el alimento.

Los procesos modernos pueden alterar la estructura de la fibra, la disponibilidad
de nutrientes y las respuestas hormonales, independientemente de la composición
final. NOVA divide los alimentos en 4 grupos distintos, no en una escala continua.

\section{Grupo 1 --- Alimentos no procesados o mínimamente procesados}

Alimentos sin transformación industrial significativa, o solo con procesos físicos
mínimos: secado, refrigeración, pasteurización.

\textbf{Ejemplos:} frutas y verduras frescas o congeladas, carnes y pescados
frescos sin aditivos, huevos, legumbres secas, cereales en grano, harina
integral, leche fresca, yogur natural sin aditivos, té, café.

\boxverde{Objetivo}{Estos alimentos deben constituir la mayor parte de cada comida.}

\section{Grupo 2 --- Ingredientes culinarios procesados}

Sustancias extraídas de alimentos del Grupo~1 o de la naturaleza, usadas en la
cocina. No se consumen solos.

\textbf{Ejemplos:} aceite de oliva, mantequilla, vinagre, sal, azúcar, miel,
harina blanca, almidón de maíz puro, jarabe de arce, leche en polvo y quesos
curados tradicionales.

Usar con moderación: el exceso de azúcar, sal y grasas saturadas sigue siendo
perjudicial aunque sean «naturales».

\section{Grupo 3 --- Alimentos procesados}

Productos del Grupo~1 con ingredientes del Grupo~2, mediante procesos como
conservación en sal, ahumado, fermentación o curación. Generalmente contienen
2--3 ingredientes.

\textbf{Ejemplos:} conservas de verduras y legumbres (con agua y sal), fruta
confitada, frutos secos tostados y salados, anchoas en aceite, jamón serrano
de calidad, quesos tradicionales (parmesano, manchego), pan artesanal, vino,
cerveza artesanal.

El Grupo~3 no hay que evitarlo; muchos alimentos culturalmente importantes
pertenecen aquí.

\section{Grupo 4 --- Ultraprocesados: el grupo a limitar}

Formulaciones industriales sin alimentos enteros o con solo una pequeña parte.
Están construidos con ingredientes extraídos, purificados o modificados
químicamente, más aditivos tecnológicos en grandes cantidades.

La señal clave: \textbf{contienen ingredientes que no existen en una cocina
doméstica.}

Ingredientes típicos: proteínas hidrolizadas, almidones modificados, jarabe de
glucosa-fructosa, aceites hidrogenados, aromas «naturidénticos», colorantes
artificiales, edulcorantes intensivos (aspartamo, sacarina, sucralosa),
emulsionantes (E471, E472), estabilizantes.

\textbf{Ejemplos:} bollería industrial, cereales azucarados, bebidas gaseosas
y energéticas, pan de molde con aditivos, salchichas reconstituidas, \textit{nuggets},
yogures aromatizados con espesantes, platos preparados congelados.

\boxrojo{Dato relevante}{En España y muchos países europeos, los ultraprocesados
aportan entre un 25 y un 35\,\% de las calorías diarias. Consumir más del 20\,\%
de las calorías diarias de NOVA~4 se asocia con riesgos para la salud.}

\section{NOVA en la compra práctica}

\textbf{1.~¿Podría cocinarlo en casa con ingredientes normales?} Si no, es
probablemente Grupo~4.

\textbf{2.~Busca las «palabras clave» del Grupo~4:} «almidón modificado»,
«proteínas de soja aisladas», «maltodextrinas», «jarabe de...», «aromas
complejos», «aceite de palma hidrogenado», «caseinato de sodio».

\textbf{3.~Regla práctica:} que más de la mitad del plato provenga del Grupo~1.

\textbf{4.~Usa la app:} E-Codes Reader muestra la clasificación NOVA de cada
producto escaneado.

\section{Tabla resumen NOVA}

{\small
\begin{tabularx}{\textwidth}{|l|L|L|l|}
\hline
\rowcolor{verdescuro}
  \textcolor{white}{\textbf{Grupo}}
  & \textcolor{white}{\textbf{Definición}}
  & \textcolor{white}{\textbf{Ejemplos}}
  & \textcolor{white}{\textbf{En la dieta}} \\
\hline
NOVA~1 & No/mín. procesados & Frutas, verduras, huevos, carne fresca & Base \\
\hline
NOVA~2 & Ingredientes culinarios & Aceite, sal, azúcar, harina, mantequilla & Moderación \\
\hline
NOVA~3 & Procesados tradicionales & Conservas, quesos, jamón, pan artesanal & Aceptable \\
\hline
NOVA~4 & Ultraprocesados & Bollería, salchichas, cereales azucarados & Limitar \\
\hline
\end{tabularx}
}

% ============================================================
\chapter{Ingredientes que vigilar en las etiquetas}

No hace falta leer cada etiqueta como un químico. Unos pocos términos clave
bastan para evitar los productos más problemáticos y tomar mejores decisiones
en segundos.

\section{Azúcares: los 10 nombres ocultos}

El azúcar se disfraza bajo múltiples nombres para parecer menos relevante en
la lista de ingredientes. Los 10 alias más comunes:

\begin{itemize}
  \item Jarabe de glucosa-fructosa
  \item Maltodextrinas
  \item Dextrosa
  \item Fructosa
  \item Zumo de caña concentrado
  \item Jarabe de maíz / jarabe de arroz
  \item Miel deshidratada
  \item Azúcar de coco
  \item Néctar de agave
\end{itemize}

\boxamarillo{La trampa del «sin azúcares añadidos»}{Un zumo «100\,\% fruta sin
azúcares añadidos» puede contener 10--12\,g de azúcares por 100\,ml, más que
algunos refrescos. La OMS recomienda mantener los azúcares libres por debajo del
10\,\% de las calorías diarias (unos 50\,g al día para un adulto de 2000\,kcal),
preferiblemente por debajo del 5\,\% (25\,g).}

\section{Grasas que evitar}

\textbf{Grasas hidrogenadas y parcialmente hidrogenadas.} Busca términos como
«aceite vegetal parcialmente hidrogenado», «grasa vegetal hidrogenada»,
«shortening». En Europa están prohibidas por encima de 2\,g/100\,g de grasa
desde 2021.

\textbf{Aceite de palma.} No es trans, pero durante el refinado a altas
temperaturas puede formar compuestos potencialmente cancerígenos (3-MCPD,
glicidol).

\boxverde{Grasas recomendadas}{Aceite de oliva virgen extra, aceite de lino,
nueces, almendras, aguacate, pescado azul. Contienen ácidos grasos
insaturados y omega-3 con efectos protectores documentados.}

\section{Aditivos controvertidos}

\textbf{Colorantes azoicos --- especial atención con niños:}
E102 (Tartrazina), E104 (Amarillo de quinoleína), E110 (Amarillo anaranjado~S),
E122 (Azorrubina), E124 (Rojo Ponceau 4R), E129 (Rojo Allura~AC).
Estudio EFSA 2007: correlación con hiperactividad en niños. En Europa es
obligatorio indicar en la etiqueta: «puede afectar a la actividad y la
atención de los niños».

\textbf{Nitratos y nitritos (E249--E252):} usados en carnes procesadas;
forman nitrosaminas potencialmente cancerígenas en presencia de proteínas
animales.

\textbf{Edulcorantes intensivos (E950, E951, E952, E954, E955):} seguros en
dosis aprobadas para adultos, pero la OMS (2023) ha expresado cautela sobre
el uso crónico para el control del peso.

\section{Sodio oculto}

Aproximadamente el 75\,\% del sodio en la alimentación occidental no se añade
durante la cocción, sino que ya está presente en los productos envasados.

\textbf{Principales fuentes:} pan, quesos, embutidos, platos preparados, salsas
en bote, caldos en pastilla, salsa de soja, ketchup, mayonesa, snacks salados,
cereales industriales.

\boxinfo{Cómo convertir sodio en sal en la etiqueta}{Multiplica el valor del
sodio por 2,5. Ejemplo: 0,6\,g de sodio = 1,5\,g de sal. Límite OMS: 5\,g
de sal al día (unos 2\,g de sodio).}

\section{La regla del 5: el filtro más rápido}

\begin{enumerate}
  \item \textbf{¿Más de 5 ingredientes?} Empieza a leer con atención.
  \item \textbf{¿Azúcares superiores a 5\,g/100\,g?} Verifica si figuran en
    los primeros puestos de la lista.
  \item \textbf{¿Sal superior a 0,5\,g/100\,g en un producto dulce?} Señal
    de ultraprocesamiento.
  \item \textbf{¿Más de 5 aditivos~E?} El producto es probablemente
    ultraprocesado.
  \item \textbf{¿No reconoces 5 o más ingredientes?} Considera dejarlo en el
    estante.
\end{enumerate}

\boxinfo{Nota}{La regla del 5 es una orientación, no una sentencia.
Úsala como punto de partida.}

% ============================================================
\chapter{Alimentos recomendados: las opciones más saludables en el supermercado}

No se trata de ser perfecto ni de comer solo ensalada. Se trata de saber qué
productos, con la misma comodidad, son genuinamente mejores que otros.

\section{El principio básico: sustituir, no eliminar}

\begin{description}
  \item[Mismo tiempo, menos procesamiento.] Cambiar yogur aromatizado por
    natural, arroz precocido por arroz integral, platos preparados por
    legumbres en conserva al natural.
  \item[El 80/20 funciona.] Si el 80\,\% de lo que comes es Grupo~1--2
    (según NOVA), el 20\,\% de productos más procesados no marca la diferencia.
  \item[Accesibilidad.] Garbanzos, huevos, harina integral, yogur natural,
    lentejas suelen costar menos que sus equivalentes procesados.
\end{description}

\section{Cereales e hidratos de carbono}

El problema no son los carbohidratos, sino su grado de refinamiento.

\textbf{Cereales recomendados:} avena en copos sin ingredientes añadidos,
arroz integral, espelta, cebada perlada, trigo sarraceno, pasta integral con
solo sémola, pan integral puro (máximo 4--5 ingredientes).

\textbf{Qué limitar:} cereales azucarados (incluso integrales con jarabe),
pan de molde con aditivos, galletas con aceite de palma, arroz precocido con
aditivos, bollería industrial.

\boxamarillo{Consejo para el desayuno}{Si los azúcares superan los 10\,g/100\,g,
estás consumiendo un producto de postre, no un desayuno saludable.}

\section{Proteínas}

\textbf{Pescado y mariscos:} pescado fresco o congelado sin rebozar, sardinas
y caballa en aceite de oliva, atún al natural, anchoas en sal. Una fuente
económica y nutritiva.

\textbf{Huevos y lácteos:} huevos enteros, yogur griego natural, requesón,
quesos frescos y curados con moderación. Evita sustitutos industriales,
yogures «proteicos» con espesantes y edulcorantes.

\textbf{Legumbres:} lentejas, garbanzos, alubias, guisantes y edamame.
Las versiones en conserva al natural son prácticas y ricas en fibra.

\section{Grasas buenas}

Las grasas de calidad ---aceite virgen extra, nueces, aguacate, pescado
graso--- son algunos de los componentes más protectores para el sistema
cardiovascular.

\textbf{Recomendadas:} aceite de oliva virgen extra, aceite de lino, nueces,
almendras, avellanas, semillas de lino y chía, aguacate, pescado azul.

\textbf{Con moderación:} mantequilla, aceite de coco, leche entera.

\textbf{Evitar:} margarinas con grasas parcialmente hidrogenadas, shortening
vegetal, aceite de palma refinado, nata vegetal con emulsionantes.

\section{Snacks inteligentes}

\begin{enumerate}
  \item \textbf{Frutos secos naturales:} nueces, almendras, avellanas sin sal
    añadida (30\,g aportan proteínas y minerales).
  \item \textbf{Fruta fresca entera:} mantiene la fibra que ralentiza la
    absorción de azúcares.
  \item \textbf{Yogur griego natural:} alto en proteínas, bajo índice
    glucémico, sin aditivos.
  \item \textbf{Queso curado + crackers integrales:} satiedad duradera con
    ingredientes simples.
\end{enumerate}

\section{Lista de la compra tipo}

{\small
\begin{tabularx}{\textwidth}{|l|L|}
\hline
\rowcolor{verdescuro}
  \textcolor{white}{\textbf{Categoría}}
  & \textcolor{white}{\textbf{Productos recomendados}} \\
\hline
Cereales & Avena, arroz integral, espelta, pasta integral, pan artesanal \\
\hline
Legumbres & Lentejas, garbanzos, alubias, guisantes \\
\hline
Proteínas & Huevos, sardinas en aceite, atún natural, yogur griego, requesón \\
\hline
Verduras & De temporada, frescas o congeladas sin aditivos \\
\hline
Grasas & Aceite de oliva virgen extra, nueces, almendras, semillas de lino \\
\hline
Fruta & Entera de temporada, frutos rojos congelados sin azúcar \\
\hline
Condimentos & Vinagre de manzana, salsa de soja poco procesada, especias \\
\hline
\end{tabularx}
}

\boxverde{La prueba del carrito}{Cuenta los productos con menos de 5 ingredientes.
Si superan la mitad del carrito, vas por buen camino.}

% ============================================================
\chapter{Por qué los alimentos ultraprocesados son malos para la salud}

No es solo cuestión de calorías o grasas: los ultraprocesados interfieren con
la saciedad, alteran el microbioma e inician inflamación crónica.

\section{Obesidad y alteración del metabolismo}

Los ultraprocesados no provocan simplemente un aumento de peso por su aporte
calórico. Están diseñados para ser hiper-apetecibles e interfieren con las
señales naturales de saciedad.

\boxinfo{Estudio Kevin Hall (NIH, 2019)}{Los participantes con dieta
ultraprocesada consumían \emph{en promedio 500 calorías más al día} que
quienes seguían una dieta de alimentos mínimamente procesados, ganando
casi un kilogramo en dos semanas.}

Los UPF contienen en promedio 2,15\,kcal por gramo ---casi el doble que los
alimentos frescos. La combinación industrial de azúcar, sal y grasas activa
los centros del placer de forma similar a ciertas sustancias adictivas.
Además, los UPF reducen el péptido~YY, la hormona intestinal que señala
saciedad al cerebro.

Estadística: \textbf{+41\,\% de riesgo de obesidad abdominal} en quienes siguen
dietas altas en ultraprocesados.

\section{Inflamación y disfunciones del sistema inmunitario}

El consumo frecuente de ultraprocesados se asocia a un aumento de las respuestas
inmunitarias inflamatorias, más allá de su simple perfil nutricional.

Sustancias habituales en los UPF ---emulsionantes, espesantes y dióxido de
titanio (E171)--- pueden alterar el microbioma intestinal y aumentar la
permeabilidad de la mucosa (\textit{leaky gut}), lo que permite que bacterias
dañinas y toxinas pasen al torrente sanguíneo y generen inflamación sistémica.

\textbf{Enfermedades autoinmunes asociadas en estudios observacionales:}
celiaquía, tiroiditis de Hashimoto, esclerosis múltiple, lupus eritematoso
sistémico, diabetes tipo~1.

\section{Salud intestinal y microbiota}

El ultraprocesamiento priva a los alimentos de fibra dietética y compuestos
vegetales bioactivos (polifenoles, antioxidantes), necesarios para nutrir
las bacterias beneficiosas del intestino.

\textbf{Disbiosis:} la reducción de la diversidad bacteriana se asocia con
resistencia a la insulina, inflamación crónica y mayor susceptibilidad a
infecciones.

\textbf{Debilitamiento protector:} los emulsionantes y edulcorantes artificiales
reducen la capa mucosa que recubre el intestino.

\textbf{Patologías más frecuentes:} enfermedad de Crohn, colitis ulcerosa,
síndrome del intestino irritable, úlceras gástricas, pólipos precancerosos
de colon.

\section{Un riesgo sistémico: múltiples mecanismos a la vez}

{\small
\begin{tabularx}{\textwidth}{|l|L|L|}
\hline
\rowcolor{verdescuro}
  \textcolor{white}{\textbf{Mecanismo}}
  & \textcolor{white}{\textbf{Efecto}}
  & \textcolor{white}{\textbf{Consecuencia}} \\
\hline
Hiper-apetecibilidad & Consumo excesivo calórico & Obesidad, diabetes tipo~2 \\
\hline
Carencia de fibra & Disbiosis & Inflamación crónica \\
\hline
Aditivos (emulsionantes) & Permeabilidad intestinal & Enf. autoinmunes, cardiovasculares \\
\hline
Azúcares/edulcorantes & Resistencia insulínica & Diabetes, síndrome metabólico \\
\hline
Sustitución de frescos & Carencia de micronutrientes & Estrés oxidativo \\
\hline
\end{tabularx}
}

Los grandes estudios epidemiológicos ---NutriNet-Santé (Francia, más de 100.000
participantes) y UK Biobank (más de 500.000 participantes)--- asocian el consumo
elevado de UPF con mayor riesgo de cáncer, enfermedades cardiovasculares,
depresión y mortalidad general.

\section{Qué hacer en la práctica}

El mensaje no es eliminar todo, sino \emph{reducir la proporción de UPF en
la dieta diaria}.

\begin{enumerate}
  \item \textbf{Identifica tus UPF habituales:} cereales azucarados, snacks,
    pan de molde, platos preparados, bebidas aromatizadas.
  \item \textbf{Sustitúyelos uno a uno:} avena en lugar de cereales azucarados,
    yogur natural en lugar de aromatizado, pan artesanal.
  \item \textbf{Aumenta la fibra --- prioridad:} legumbres, verduras, cereales
    integrales, fruta entera.
  \item \textbf{Usa apps para identificar los ocultos:} muchos productos
    aparentemente saludables son NOVA~4.
  \item \textbf{No busques perfección:} el objetivo es que la mayor parte de
    lo que comas cada semana sea comida real.
\end{enumerate}

\boxverde{Buenas noticias}{Incluso una mejora del 20--30\,\% en la calidad
de la dieta produce efectos medibles en la salud a medio plazo.}

% ============================================================
\chapter{Seguridad alimentaria para niños y ancianos}

Los menores y adultos mayores reaccionan de forma diferente a los aditivos
e ingredientes de baja calidad. Esta guía explica qué revisar antes de
comprar para estos grupos vulnerables.

\section{Por qué son más vulnerables}

\textbf{Niños: bajo peso corporal, sistemas inmaduros.} Un niño de 15\,kg que
consume una merienda de 30\,g con colorantes recibe una dosis relativa 4--5
veces superior a la de un adulto de 70\,kg. El hígado y los riñones
infantiles metabolizan ciertas sustancias más lentamente.

\textbf{Mayores: función renal y hepática reducidas.} La capacidad de
eliminar ciertas sustancias disminuye con la edad. La polifarmacia puede
interactuar con aditivos. El sodio excesivo agrava la hipertensión y la
insuficiencia renal.

\textbf{Efectos acumulativos:} un producto ocasional no causa problema, pero
el consumo diario de colorantes en múltiples comidas genera efectos relevantes.

\section{Aditivos críticos para niños}

\textbf{Colorantes azoicos (a evitar):}
E102, E104, E110, E122, E124, E129. En Europa deben indicar en la etiqueta:
«puede afectar a la actividad y la atención de los niños». Se encuentran en
caramelos de colores, bebidas naranjas/rojas, gelatinas, bollería y
cereales de desayuno.

\textbf{Edulcorantes intensivos (no autorizados para lactantes y pequeños):}
E951, E950, E955, E954. Presentes en yogures «0\,\% azúcar», bebidas
\textit{light}, jarabes y caramelos sin azúcar.

\boxrojo{Cafeína}{Una lata de 250\,ml de bebida energética puede contener
hasta 80\,mg de cafeína ---casi el doble del límite recomendado para un
niño de 15\,kg.}

\section{Qué limitar en personas mayores}

\textbf{Sodio:} principal riesgo para la hipertensión y la insuficiencia renal.
Objetivo: menos de 5\,g de sal diarios. Principales fuentes: platos preparados,
embutidos, quesos curados, caldos industriales.

\textbf{Fosfatos (E338--E452):} interfieren con la absorción del calcio,
un problema para personas con riesgo de osteoporosis. Comunes en productos
cárnicos procesados y quesos fundidos.

\textbf{Interacciones con medicamentos:} la vitamina~K de las verduras de
hoja verde interactúa con anticoagulantes; el pomelo con estatinas; el zumo
de arándanos con ciertos anticoagulantes.

\section{Productos «para niños» que conviene revisar}

\textbf{Cereales de desayuno «para niños»:} suelen contener 25--35\,g de
azúcares por 100\,g, colorantes artificiales y vitaminas añadidas para
compensar la baja calidad nutritiva de la base.

\textbf{Bebidas «a base de zumo»:} pueden contener solo el 10--30\,\% de zumo
real; el resto es agua, azúcares, acidulantes y colorantes. Verifica siempre
el porcentaje de fruta en la etiqueta.

\textbf{Yogures «para niños» con frutas:} contienen más azúcar que las versiones
para adultos, más espesantes y aromas artificiales. El yogur griego natural
con fruta fresca añadida es una alternativa mejor.

\section{Lista segura para la compra}

\begin{enumerate}
  \item \textbf{Para niños:} elige productos con pocos ingredientes reconocibles;
    evita los que contengan más de 2--3 colorantes o edulcorantes artificiales.
  \item \textbf{Verifica los azúcares en los «productos para niños»:}
    menos de 10\,g/100\,g en desayunos y snacks.
  \item \textbf{Para mayores:} prioridad a bajo sodio; menos de 0,3\,g de
    sal/100\,g.
  \item \textbf{Escanea con la app antes de comprar:} E-Codes Reader indica
    de inmediato colorantes azoicos, edulcorantes inadecuados y nivel de sodio.
  \item \textbf{Prioriza el Grupo~1--2 (NOVA) para las comidas principales:}
    verduras, legumbres, huevos, pescado, yogur natural, fruta entera.
\end{enumerate}

\section{Tabla comparativa}

{\small
\begin{tabularx}{\textwidth}{|L|l|l|}
\hline
\rowcolor{verdescuro}
  \textcolor{white}{\textbf{Aditivo/ingrediente}}
  & \textcolor{white}{\textbf{Niños}}
  & \textcolor{white}{\textbf{Mayores}} \\
\hline
Colorantes azoicos (E102, E110, E122\ldots) & Evitar & Limitar \\
\hline
Edulcorantes intensivos & No para pequeños & Con moderación \\
\hline
Sodio elevado (>1,5\,g sal/100\,g) & Limitar & Evitar \\
\hline
Fosfatos (E338--E452) & Moderación & Limitar \\
\hline
Cafeína & Evitar <12 años & Atención \\
\hline
Nitratos/nitritos (E249--E252) & Limitar & Limitar \\
\hline
\end{tabularx}
}

% ============================================================
\backmatter

\chapter*{Créditos y descarga de la app}
\addcontentsline{toc}{chapter}{Créditos}

\begin{center}
\vspace{1cm}
{\large\bfseries\color{verdescuro} Come Consciente}\\[0.4cm]
{\itshape Guía práctica sobre aditivos alimentarios}\\[0.3cm]
{\small Con E-Codes Reader para Android}

\vspace{1.5cm}
\textbf{Versión 1.0} \quad \textbullet \quad 2026\\
\textbf{Licencia:} Creative Commons BY-NC 4.0\\[0.3cm]
\textbf{Sitio web:} \href{https://faustobe.it}{faustobe.it}\\
\textbf{Aplicación:} \href{https://play.google.com/store/apps/details?id=it.faustobe.ecodes}{E-Codes Reader en Google Play}

\vspace{1.5cm}
{\small Escanea el código QR para descargar la app:}\\[0.5cm]

\qrcode[height=4cm]{https://play.google.com/store/apps/details?id=it.faustobe.ecodes}

\vspace{1cm}
{\footnotesize Este ebook fue producido en \LaTeX{} por faustobe.}
\end{center}

\end{document}
